\documentclass[conference]{IEEEtran}
\usepackage[T1]{fontenc}
\usepackage[utf8]{inputenc}
\usepackage{lmodern}
\usepackage{amsmath,amssymb}
\usepackage{xcolor}
\usepackage{booktabs}
\usepackage{enumitem}
\usepackage[numbers,sort&compress]{natbib}
\usepackage{hyperref}
\hypersetup{colorlinks=true,linkcolor=blue,citecolor=blue,urlcolor=blue}

% --- draft machinery: remove before camera-ready ---
\newcommand{\placeholder}[1]{\textcolor{red}{\textbf{[PLACEHOLDER: #1]}}}
\newcommand{\wj}[1]{\textcolor{blue}{\emph{[WJ: #1]}}}

\begin{document}

\title{The Certified-Human Anchor: Temporal-Holdout Validation\\ for Corpus-Level AI-Content Estimators}

\author{\IEEEauthorblockN{Anonymous Author(s)}
\IEEEauthorblockA{Anonymous Institution(s)\\anonymous@example.org}}

\maketitle

\begin{abstract}
Corpus-level AI-content estimators report what fraction of a corpus is AI-generated by fitting a mixture against a ``certified-human'' reference, almost always drawn from text predating ChatGPT's late-2022 release on the assumption that such text cannot be AI-contaminated. We identify the \emph{purity--validity gap} in this convention: a pre-2022 anchor can be label-perfect and still license a false error guarantee, because the false-positive guarantee belongs to the interaction between a detector, a threshold, and the distribution of human text --- not to the labels --- and human writing distributions move over the years a ``pre-2022'' pool spans. We measure the gap where a user of the estimator feels it: holding a target corpus fixed and swapping only the age of the human reference, we find the estimated AI-fraction shifts by \placeholder{$\approx$ZZ~pp across an 8--14-year reference span (pre-registered floor: 2~pp)}. Behind that shift, a temporal-holdout protocol over year-stratified, hard-timestamped corpora and three pre-registered detector families measures a false-positive drift of \placeholder{$\approx$XX~pp per decade}, attributed to genuine stylistic change only after surviving leakage-immune, composition-matched, and pipeline-constant controls, and bounded net of a human-audited anchor-impurity estimate. The output is a validity envelope --- time, genre, domain --- replacing a field-wide assumption with a reported, falsifiable quantity; both the drift branch and the estimator-self-robustness branch are pre-registered as publishable findings.
\end{abstract}

\begin{IEEEkeywords}
AI-generated text detection, corpus-level prevalence estimation, temporal validity, measurement infrastructure, pre-registration, empirical software engineering
\end{IEEEkeywords}

\section{Introduction}\label{sec:intro}

Estimating how much of a corpus is AI-generated has increasingly shifted from a per-document classification problem to a population-level measurement problem. Rather than asking whether any single passage was written by a model, a growing line of work asks what fraction of an entire corpus --- a peer-review pool, a public forum, a submission queue --- was touched by one, and answers it by fitting a mixture model against two reference distributions: an AI-generated reference and a ``certified-human'' reference \cite{MonitorAIScale}. That human reference is almost always assembled from text dated with confidence to before ChatGPT's public release in November 2022, on the reasoning that text written before large language models were widely accessible cannot itself be LLM-contaminated. This pre-2022 anchor has become a load-bearing but largely unexamined convention: nearly every corpus-level estimator that needs a human reference reaches for it, yet its validity is asserted by community practice, not measured by any tool that depends on it.

We call the resulting blind spot the \textbf{purity--validity gap}: a pre-2022 anchor's authorship purity and its statistical validity are different properties, and they come apart under time. A document dated 2015 can be genuinely, unambiguously human-written and still be flagged as AI at a rate that has drifted upward relative to whatever rate the anchor was calibrated to provide --- for reasons that have nothing to do with any label being wrong. The false-positive guarantee does not belong to the document label; it belongs to the interaction between a detector, a decision threshold, and the distribution of human text being evaluated. Human writing style is not static, so if it shifts across the calendar years a ``pre-2022'' pool spans, an estimator calibrated against recent pre-2022 text will flag older pre-2022 text at an elevated rate, inflating reported AI-prevalence even though every document in the pool is exactly as certified-human as the convention assumes. The guarantee can decay while every label stays correct.

This gap sits directly downstream of an active research line, yet no prior work measures it. The mixture-estimation framework that popularized the pre-2022 anchor \cite{MonitorAIScale} fits its mixing coefficient against a fixed human reference and treats it as unproblematic ground truth; its own text flags temporal token-frequency shift as an added source of error but never measures it --- an instrument this paper audits directly rather than reuses. A longitudinal concept-drift benchmark in a different domain \cite{LAMDA} supplies the year-stratified validation template this paper adapts, but targets classification-accuracy decay in malware detection, not a false-positive bound for a text negative-anchor. And a cluster of work on contamination- and membership-detector fragility \cite{BenchContamFragility,MinKpp,ContamICL,MarkovCalib} shows that ``seen-versus-unseen'' separability collapses as compared distributions converge or shift, but stops at the diagnosis without turning it into an explicit bound as a function of temporal distance. A systematic search of this literature (11,683 records, five angles) returns no work that measures the pre-2022 anchor's temporal validity --- the gap is genuinely open, not merely unfashionable.

Turning that gap into a measurable claim is hard for three reasons a bare false-positive-rate curve cannot resolve. An observed rise over calendar years is equally consistent with genuine stylistic drift, with ordinary corpus-composition confound (older text skewing toward different genres, topics, or lengths), or with undetected AI content leaking into nominally pre-2022 buckets --- the curve alone separates none of the three. The ``certified-human'' label is itself a metadata proxy rather than a verified property, so any measured drift carries a residual, unresolvable uncertainty unless that uncertainty is explicitly bounded. And what a downstream user consumes is not a per-document statistic but a corpus-level prevalence number, whose mixing coefficient is a nonlinear function of the entire reference distribution --- so a nonzero document-level drift does not, by itself, establish that anything a practitioner reads off an estimator actually moves.

We turn each obstacle into a specific design commitment rather than a hedge. A leakage-immune restriction to years unambiguously predating large-scale LLM availability and a composition-matched resampling that equalizes topic and length across year buckets each rule out exactly one rival explanation, while a synthetic-drift injection independently confirms that style-like perturbation alone can move the measured statistic in the observed direction. The anchor's own residual impurity is bounded through a human-audited, inter-rater-validated spot-check rather than assumed to be zero. Whether any of this reaches the number a practitioner reports is tested directly --- not inferred --- by holding a target corpus fixed and swapping only the age of the human reference fed into \cite{MonitorAIScale}'s own mixture estimator. And the resulting validity envelope is drawn from a confidence bound rather than a raw significance threshold, so that ``not enough data to tell'' and ``genuinely safe'' are never reported as the same finding.

We evaluate this design with a fixed-target, swapped-reference protocol at its center: a target corpus is held constant while only the age of the human reference fed into \cite{MonitorAIScale}'s mixture estimator changes, so that any movement in the reported AI-fraction is attributable to anchor age alone. Around that core sit year-stratified, hard-timestamped corpora spanning more than a decade of pre-2022 history --- including an SE-native year-stratified corpus (a venue's own pre-2022 review/abstract archive), since the instrument under audit was born in the peer-review setting --- a pre-registered battery of three named detector families (\cite{DetectGPT} curvature, \cite{LearnToDistance} rewrite-likelihood, \cite{DnR} black-box recovery; families named before any data, so the $\geq$2-of-3 replication rule is confirmatory, and any later-added family is exploratory and separately reported), and a year-shuffle placebo plus a zero-distance sanity cell guarding the pipeline against manufacturing --- or masking --- its own trend. In summary, we contribute:

\begin{itemize}
\item \textbf{A propagation measurement --- what moves, and by how much:} holding a target corpus fixed and swapping only the human reference's age in \cite{MonitorAIScale}'s estimator, we measure \textbf{[PLACEHOLDER: a $\approx$ZZ-pp shift in estimated AI-fraction across an 8--14-year reference span, against the pre-registered $\geq$2pp floor]} --- the bridge from a document-level statistic to the corpus-level number a venue actually reports and acts on. Both branches are pre-registered: propagation at or above the floor indicts the convention; invariance below it delivers the reassuring co-headline that the estimator is self-robust to anchor age even where documents drift.
\item \textbf{A new estimand and its measurement --- why it moves:} the false-positive rate of a corpus-level AI-content estimator on certified-human text as a function of temporal distance from the anchor's certification date --- a quantity no prior estimator defines --- measured at \textbf{[PLACEHOLDER: $\approx$XX percentage points per decade, against the pre-registered $\geq$5\%/year floor]} and yielding a validity envelope that \textbf{[PLACEHOLDER: closes after $\approx$YY years]} at 95\% confidence, with an operational thin-cell rule (buckets below $n_{\min}=1{,}000$ or above 1 pp CI half-width report INCONCLUSIVE, never ``safe'').
\item \textbf{An attribution design --- what it is:} four controls, each immune to exactly one rival explanation --- a leakage-immune pre-2019 sub-range (with a $\leq$2016 sensitivity cut), a composition-matched resample over named covariates (topic, length, genre/register, source), a provenance-stratified pipeline-constant arm that separates the history of writing from the history of text processing, and a synthetic-drift injection --- with the drift attributed to genuine stylistic change only where it \textbf{[PLACEHOLDER: survives the leakage-immune, composition-matched, and pipeline-constant arms for $\geq$2 of 3 detector families]}.
\item \textbf{The purity--validity gap itself, named and formalized:} anchor purity and anchor validity are properties of different objects, which is what lets a label-perfect anchor license a decaying guarantee --- the identification the three measurements above operationalize into a reusable temporal-holdout validation protocol, and a measurement-infrastructure validity contribution: any research community that reuses prevalence estimators inherits the anchor, and this protocol is how it checks what it inherited.
\end{itemize}

\section{Motivation}\label{sec:motivation}

The pre-2022 anchor is not a peripheral calibration nuisance; it is the single most load-bearing ingredient in an emerging measurement paradigm. As detection research has shifted from unreliable per-document classification toward corpus-level prevalence estimates \cite{MonitorAIScale}, the human reference distribution used to define ``prevalence'' is no longer one input among several --- it is the object every subsequent number in the pipeline is computed relative to. If that reference's validity erodes as the anchor ages further past 2022, every downstream estimate that reuses it inherits an unstated, directional bias, and does so silently: a prevalence number that looks precise while resting on an unmeasured temporal assumption.

The consequence is a specific and dangerous ambiguity. Suppose a venue reports, via \cite{MonitorAIScale}-style estimation against a fixed 2013--2015 human reference, that the AI-involved fraction of its submissions rises year over year. If this paper's central hypothesis holds, part of that reported rise need reflect no change in actual AI use at all --- it can be an artifact of the reference growing one year more temporally distant from each successive cohort. From outside the pipeline, a genuine increase in AI-assisted writing and a compounding anchor-age artifact are observationally identical: both produce the same upward-trending number, and only an explicit, measured validity envelope can tell a reader which one they are seeing.

What makes the gap easy to overlook is that the pre-2022 anchor looks like one of the few assumptions in AI-content estimation that is trivially certifiable. A timestamp seems to hand over a clean negative label for free. But certifying that a document predates widespread LLM use --- the easy part, a matter of metadata --- is not the same as certifying that its distribution still supplies the error guarantee the estimator was calibrated to expect --- the hard part, and the one no timestamp settles. This paper's premise is that the second certificate has to be measured, not assumed to follow from the first.

That is also why the measurement must be year-indexed rather than a single pooled audit of whether pre-2022 text is ``safe on average.'' A pooled check can only report whether the anchor is uniformly and mildly noisy across its span; it cannot distinguish that from an anchor safe for recent years but already unsafe for the oldest ones --- the difference between a flat guarantee and a validity envelope with a real, locatable boundary. Only a year-by-year reading finds where that boundary sits, which is exactly what a downstream researcher needs before deciding how far back into the pre-2022 pool a human reference can safely be drawn.

The stakes extend past any single estimator. Researchers building or reusing corpus-level estimators need to know whether the calibration convention stays safe as models and corpora age further past 2022. Venues, platforms, and institutions that cite prevalence numbers as evidence for policy or enforcement decisions are relying, usually without stating it, on an anchor whose temporal shelf life has never been reported. And the broader contamination-detection literature, which treats ``human-versus-AI'' gaps as stable analytic devices, has not asked whether its most common such gap has an expiration date.

Every part of this study is pre-registered against a single binary, falsifiable compass criterion before any data is collected (our pre-registration protocol). The design is therefore decision-relevant whichever way it resolves: a genuine, propagating drift means every reuse of the anchor without a stated envelope carries a directional bias that grows as the anchor ages; a flat, non-propagating result instead delivers the first quantified license --- bounded by an explicit envelope --- for a convention the field has so far only assumed.

\section{Problem Definition}\label{sec:problem}

\noindent\textbf{Formal setup.} Let $y^*$ denote the certification cutoff (2022) and let $\Delta(y) = y^* - y$ denote the temporal distance of a calendar-year bucket y. A pre-2022 anchor corpus is partitioned into year-buckets $\{A_y\}$, each containing hard-timestamped, certified human-authored documents drawn from a fixed source. For an estimator $D_k$ operating at a fixed decision threshold, let $\mathrm{FPR}_k(y)$ denote its empirical false-positive rate on $A_y$ --- the fraction of certified-human documents in bucket y that $D_k$ flags as AI-generated. The object of study is $\mathrm{FPR}_k(y)$ as a function of $\Delta(y)$: does it stay flat (the anchor's guarantee is uniform across the pre-2022 span), or does it grow with temporal distance (the guarantee erodes for older text)? The \textbf{validity envelope} $E_k(\epsilon, \delta)$ is the largest contiguous year-range whose upper confidence bound on $\mathrm{FPR}_k(y)$ stays within a pre-declared tolerance $\epsilon$ of the nominal calibrated rate, at confidence level $1-\delta$ --- the interval construction is standard Clopper--Pearson; the contribution is the estimand it bounds, a temporal negative-anchor validity no prior estimator defines. A bucket enters the envelope only if $n_y \geq n_{\min} = 1{,}000$ and its realized confidence half-width is $\leq$ 1 pp; a bucket failing either trigger reports \textbf{INCONCLUSIVE}, never ``safe'' --- the operational rule behind the principle that not enough data is not the same thing as no drift. A second object of study is the bias this induces in a downstream mixture-based prevalence estimator $\hat{\pi}(T; H)$, which estimates the AI-fraction of a fixed target corpus T given a human reference H: does swapping H across year-buckets shift $\hat{\pi}(T; H)$ --- that is, does drift measured at the document level actually move the number a downstream user consumes?

\smallskip\noindent\textbf{Why this is hard.} Any relationship between $\mathrm{FPR}_k(y)$ and $\Delta(y)$ is compatible with at least four different mechanisms --- genuine secular change in human writing style, a corpus-composition confound (older buckets skewing toward different genres, topics, or lengths than newer ones), undetected AI-content leakage into nominally pre-cutoff buckets, and a *processing-history* confound: OCR error profiles, encoding artifacts, and house-style normalization correlate with the calendar year of digitization rather than the year of authorship, and shift token distributions in exactly the direction a token-frequency estimator reads as ``less like the reference'' --- and the curve alone cannot distinguish them. Before any of these can even be evaluated, the composition and provenance confounds must be ruled out: if y is entangled with corpus composition or with the digitization pipeline, a change in $\mathrm{FPR}_k(y)$ need not be evidence about writing at all. Compounding this, the label ``certified human'' that defines each $A_y$ is a metadata proxy --- a publication-year timestamp --- rather than a verified property of authorship, so any residual mislabeling contributes an unresolvable source of noise unless it is explicitly bounded. Finally, $\mathrm{FPR}_k(y)$ is a document-level statistic, while $\hat{\pi}(T; H)$ is a nonlinear function of the entire reference distribution H; a nonzero drift in $\mathrm{FPR}_k(y)$ does not by itself imply any particular shift in $\hat{\pi}(T; H)$, so the practical stakes of the drift can only be established by intervening on H directly, not by extrapolating from the document-level curve.

\section{Study Design}\label{sec:design}

At a conceptual level, the study is a single audit pipeline with five cooperating parts that turn a suspected but unmeasured anchor weakness into a bounded, attributed, and propagated quantity. Year-stratified anchor text and an existing battery of detectors and a mixture-based prevalence estimator enter the pipeline; a pair of disentanglement controls turns a raw false-positive-rate trend into an attributed claim about which of three rival mechanisms produced it; a purity-bounding step turns the anchor's own unverifiable authorship claim into an explicit uncertainty rather than a silent assumption; a reference-swap procedure turns a document-level drift signal into a measured shift in the corpus-level number a downstream estimator actually reports; and a confidence-bound envelope construction turns the fitted drift curve into a stated boundary condition rather than an all-or-nothing verdict.

The \textbf{disentanglement controls} take the raw false-positive-rate series for each year bucket and each detector and test it against two independent, mechanism-specific alternatives --- a leakage-immune sub-range and a composition-matched resample --- plus a synthetic-drift injection that checks whether style-like perturbation alone can reproduce the observed direction. The \textbf{purity-bounding procedure} takes a small human-audited sample of the anchor sources and turns it into an explicit residual-impurity estimate that is reported alongside, not folded silently into, every downstream bound. The \textbf{composition-matching procedure} takes the raw year-bucket samples and produces a matched resample in which topic and length distributions no longer vary with year, isolating whether the drift signal is genuinely temporal. The \textbf{reference-swap propagation test} takes a fixed target corpus and the mixture-based prevalence estimator and measures how much the estimated AI-fraction moves when only the age of the human reference changes. The \textbf{confidence-bound envelope construction} takes the fitted drift curve and produces the largest contiguous year-range whose upper confidence bound on false-positive rate stays within a pre-declared tolerance of the nominal calibrated rate, distinguishing a genuinely certified boundary from one that is merely under-sampled.

\subsection{Disentanglement Controls}

\textbf{Challenge.} An observed false-positive-rate slope over calendar years is compatible with three different mechanisms --- genuine stylistic drift, undetected AI-content leakage into nominally pre-2022 buckets, or staleness of the estimator itself --- and no single statistic separates them.

\textbf{Insight.} No one control rules out all three at once, but each control can be built to be blind to exactly one alternative, so requiring survival across independently-designed controls narrows the explanation by elimination rather than by assumption.

\textbf{Method.} The drift regression is refit restricted to years unambiguously predating large-scale LLM availability (a leakage-immune sub-range), where any positive slope cannot be explained by leaked AI content by construction. The sub-range boundary is set at 2018 because GPT-2's February 2019 release marks the first widely available model producing fluent open-domain text at scale, so hard-timestamped text at or before 2018 predates any plausible LLM-generated leakage; a $\leq$2016 sensitivity cut is reported alongside, and the two cuts must agree in sign for the leakage-immunity claim to stand, so the conclusion is visibly boundary-insensitive. Separately, graded lexical and syntactic perturbations are injected into a single fixed year's human sample at increasing intensity, and the flag rate is checked for a monotone rise with intensity --- establishing, independently of any real year-over-year data, that style-like change alone can move the measured statistic in the observed direction. A slope's sign is required to replicate across a majority of independently-built detector families before it is attributed to genuine drift rather than any one estimator's idiosyncrasy.

\subsection{Anchor Purity Bounding}

\textbf{Challenge.} ``Certified human'' is a metadata proxy --- a publication-year timestamp --- not a verified property of authorship, and no scalable oracle exists to certify that a nominally pre-2022 document involved no later revision, ghostwriting, or mislabeling.

\textbf{Insight.} Purity cannot be certified exactly at corpus scale, but the uncertainty it contributes can be bounded and folded into the reported result instead of silently assumed to be zero --- the same bound-rather-than-assume move this paper applies to the anchor's temporal validity, applied recursively to the anchor's own construction.

\textbf{Method.} Anchor construction is restricted to sources carrying hard, third-party archival timestamps rather than self-reported metadata. A human-audited spot-check of $n=2{,}000$ documents, stratified proportionally across year buckets and sources, is then reviewed for signs of misdating, later heavy revision, translation, or ghostwriting --- the sample size is powered against the envelope's tolerance: at a plausible true impurity of 0.3\%, $n=2{,}000$ yields a one-sided exact-binomial 95\% upper bound of roughly 0.7\%, below the 1\% floor the envelope math requires, so the audit can demonstrate compliance rather than merely fail to reject it. Two independent annotators label the full sample against a rubric piloted on ~100 documents, with inter-rater agreement of $\kappa \geq 0.7$ required before the estimate is trusted and disagreements resolved by a third annotator's majority vote. The resulting impurity upper bound enters the analysis in exactly one place: it widens each year-bucket's false-positive confidence bound additively before the validity envelope is drawn, so impurity uncertainty can push a bucket to INCONCLUSIVE but can never manufacture a spurious ``safe.''

\subsection{Composition-Matched Resampling}

\textbf{Challenge.} Calendar year and corpus composition --- genre, topic, document length --- co-vary in almost every available continuously-archived corpus, so an apparent false-positive-rate trend over years may reflect composition rather than time.

\textbf{Insight.} The only way to know whether year itself is doing causal work is to build a comparison in which composition cannot vary, and check whether the trend survives when that variation is removed.

\textbf{Method.} A single continuously-archived source is fixed per experiment so that year is the only structural variable, with per-year token budgets and sample sizes matched across buckets. Each year bucket is additionally resampled to equalize its distribution over the exhaustively named matched covariates --- topic cluster, document length, genre/register, and source --- against a reference bucket, and the drift regression is refit on this composition-matched resample. The observed slope is attributed to genuine temporal change only if it survives --- same sign, comparable magnitude, still significant --- in this matched arm.

\subsection{Provenance-Stratified (Pipeline-Constant) Control}

\textbf{Challenge.} Composition matching covers what a document is about and how long it is, but not how it was digitized: OCR error profiles, encoding artifacts, house-style normalization, and markup-stripping differences correlate with the calendar year of digitization rather than the year of authorship, and they shift token distributions in exactly the direction a token-frequency estimator reads as ``less like the reference'' --- a drift curve could be measuring the history of text processing, not the history of writing.

\textbf{Insight.} The processing confound cannot be matched away statistically, because it is a property of the corpus assembly pipeline rather than of the documents --- but it can be held constant by construction, by restricting one arm to text that never passed through a heterogeneous pipeline at all.

\textbf{Method.} Within at least one corpus, the digitization pipeline is held constant across all years: born-digital text only, from a single source, ingested through a single processing pipeline (for example, arXiv abstracts retrieved from the arXiv's own metadata API, or a single venue's own review archive). The year slope is re-estimated on this pipeline-constant arm, and the stylistic-drift attribution stands only if the slope survives there with the same sign and pre-registered significance. If the slope vanishes in the pipeline-constant arm, the drift is re-attributed to processing artifacts --- itself a publishable finding with a pre-committed wording: the anchor is unsafe as conventionally assembled from mixed-provenance pools, but safe when restricted to pipeline-constant sources, making the validity envelope a provenance condition rather than a time condition.

\subsection{Reference-Age Propagation Test}

\textbf{Challenge.} A per-document false-positive-rate drift is not automatically informative about the object a downstream user actually consumes --- a corpus-level prevalence number --- because the mixture estimator's mixing coefficient is a nonlinear function of the whole reference distribution, not a simple pass-through of document-level error.

\textbf{Insight.} The only way to measure how much anchor age matters to the reported number is to intervene on the reference directly and observe the corpus-level estimate move, holding everything else about the target corpus fixed.

\textbf{Method.} A target corpus is held fixed --- both a synthetic mixture of known AI-fraction and a real post-2022 corpus of unknown ground truth --- while the human reference bucket fed into the mixture-based prevalence estimator is swapped between a recent bucket and the most temporally distant tested bucket. The resulting shift in the estimated AI-fraction is read off directly as the measured propagation of anchor-age drift into the corpus-level estimand, rather than inferred analytically from the document-level slope.

\subsection{Confidence-Bound Validity Envelope}

\textbf{Challenge.} Per-year sample sizes shrink for older buckets, widening confidence intervals, so drawing a validity boundary at the first year a raw significance test crosses a threshold risks confusing under-power with genuine safety.

\textbf{Insight.} A rate estimate's confidence-interval width is itself information about whether a boundary is certified or merely unmeasured, so building the envelope's edge from the confidence bound rather than the point estimate keeps these two failure modes from being silently conflated.

\textbf{Method.} The validity envelope is defined as the largest contiguous year-range whose upper confidence bound on false-positive rate stays within a pre-declared tolerance of the nominal calibrated rate, at a stated confidence level. This construction is compared directly against a naive envelope drawn at the first year a raw significance test crosses its threshold, to check that the confidence-bound construction is doing genuine work rather than agreeing with the simpler rule by coincidence.


\section{Experiment Design}\label{sec:experiments}

\subsection{Experimental Setup}

\textbf{Datasets.}  arXiv abstracts, 2007--2022, drawn from a single continuously-archived source --- the primary pilot source, whose hard timestamps satisfy the purity and leakage constraints directly, whose $\geq$15-year span provides the pre-2018 years needed for the leakage-immune sub-range control, and whose born-digital, single-pipeline provenance makes it the natural pipeline-constant arm for RQ8. An SE-native year-stratified corpus --- a venue's own pre-2022 peer-review or abstract archive (e.g., ICLR/NeurIPS reviews 2013--2021, the setting in which the audited mixture estimator \cite{MonitorAIScale} was originally built) --- anchoring the study in the community whose measurement infrastructure it validates: an SE or ML venue deciding whether its reported AI-prevalence trend is real relies on exactly this envelope. CC-News, a dated news archive --- a scale-up genre with faster topical churn, used to stress-test whether the composition-matching procedure and the drift result generalize beyond formal academic prose. Project Gutenberg, dated by first publication --- extends the anchor's temporal depth well beyond what any recent corpus offers, testing whether the confidence-bound envelope genuinely narrows for very old buckets or merely reflects thinner samples; its heterogeneous digitization history also makes it the maximal-contrast comparison for the pipeline-constant arm. Reddit/Pushshift comments, post-date stratified --- an informal, high-topic-churn genre used as an additional composition-matching stress test and to diversify the genre coverage the validity envelope must eventually bound. HC3 (Human-ChatGPT Comparison Corpus) --- supplies the AI-side reference needed to build a known-ground-truth mixed target corpus for the reference-age propagation test, where the true AI fraction is known by construction. RAID (Robust AI Detection benchmark) \cite{RAID} --- used exclusively as a pre-registered instrument-validity check: before any threshold is frozen on the calibration split, each instance-level detector is confirmed to reproduce its expected behavior on RAID's known-clean human/AI samples, underwriting the no-threshold-leakage assumption every research question below depends on; it is not itself a claim-bearing dataset for any result.

\smallskip\noindent\textbf{Detector battery and estimator under audit.} Three instance-level detector families, named before any data so the $\geq$2-of-3 replication rule is confirmatory: a zero-shot probability-curvature detector \cite{DetectGPT}, a rewrite-likelihood detector \cite{LearnToDistance}, and a single-call black-box recovery detector \cite{DnR}. The mixture-based prevalence estimator \cite{MonitorAIScale} is not a battery member: it is the propagation object whose human reference is swapped in RQ6. Any detector family added after data collection begins is exploratory and separately reported.

\smallskip\noindent\textbf{Baselines.}  (validity-gate and sanity-check comparators, not competitors to beat --- this is a measurement study, and no research question here is framed as ``our procedure beats baseline X''). An anchor-as-asserted control reports current practice directly: assume the nominal calibrated false-positive rate holds uniformly across the whole pre-2022 span; the paper's premise only has stakes if this control is actually violated. A year-shuffle placebo permutes year labels within a source and refits the drift model --- the placebo slope must be null, or the measurement pipeline itself is manufacturing an artifact, a validity gate on the instrument rather than a competitor. Its companion positive control is a zero-distance sanity cell: the detector evaluated on a held-out recent slice disjoint from the calibration split (temporal distance $\approx$ 0) must report a false-positive rate within the confidence interval of the nominal calibrated rate, demonstrating the pipeline can report ``no drift'' where none should exist --- the placebo shows the pipeline cannot invent a trend; the sanity cell shows it does not mask one either. The composition-matched resampled arm and the pipeline-constant provenance arm are the load-bearing rival-explanation controls described in the Methodology, not competitors to beat.

\smallskip\noindent\textbf{Metrics.}  Per-bucket false-positive rate and its confidence interval; the drift regression's slope and significance, with and without each disentanglement control applied; the estimated residual anchor-impurity rate; the shift in estimated AI-fraction under reference-year swaps; the boundary year of the validity envelope under both the confidence-bound and naive constructions.

\smallskip\noindent\textbf{Pre-registered assumptions and gates.} All thresholds, floors, and analysis gates below were fixed before any data collection; the full pre-registration document accompanies the replication package. \textbf{(A1) Anchor purity.} Hard-timestamped pre-cutoff sources contain $<1\%$ mis-provenanced text, verified by the frozen $n=2{,}000$ two-annotator audit (RQ4); the impurity upper bound widens per-bucket confidence intervals before the envelope is drawn. \textbf{(A2) Timestamp accuracy.} Publication-year metadata errors are rare and non-systematic, cross-checked on a 200-document sample against independent archival snapshots. \textbf{(A3) Within-source comparability.} Year-buckets within one fixed source are comparable after covariate matching; violated if and only if the composition-matched arm (RQ5) kills the slope. \textbf{(A4) No threshold leakage.} Every detector threshold is frozen once on a held-out 2021 calibration split, disjoint from all evaluation buckets, before any evaluation false-positive rate is computed; thresholds are hashed and logged. \textbf{(A5) Estimator reproducibility.} The mixture prevalence estimator can be re-run with a swapped human reference; if irreproducible, a re-implemented estimator of the same family is substituted and declared. \textbf{(A6) Power and the thin-cell rule.} With $\approx$3{,}000 documents per bucket, a 1-pp absolute FPR difference is detectable with power $\geq 0.9$ at the Bonferroni-corrected level; a bucket enters the envelope only if $n_y \geq 1{,}000$ and its Clopper--Pearson half-width is $\leq 1$~pp, else it reports inconclusive.

\smallskip\noindent\textbf{Research questions.}
\begin{enumerate}[leftmargin=1.6em]
\item RQ1 --- Does the pre-2022 anchor license a bounded, non-uniform false-positive guarantee at all, as opposed to a flat one or none?
\item RQ2 --- Does the false-positive-rate drift survive restriction to a leakage-immune, pre-LLM sub-range of years?
\item RQ3 --- Does synthetic stylistic perturbation alone reproduce the direction of the observed drift?
\item RQ4 --- How much residual impurity does the certified-human anchor actually carry?
\item RQ5 --- Does the false-positive-rate drift survive composition-matched resampling?
\item RQ6 --- Does anchor age bias the corpus-level prevalence estimate a downstream user actually consumes?
\item RQ7 --- Does a confidence-bound validity envelope draw a meaningfully different boundary than a naive significance-threshold envelope?
\item RQ8 --- Does the false-positive-rate drift survive a provenance-stratified, pipeline-constant arm? (Attribution rank: alongside RQ2 and RQ5 --- the third one-alternative-immune control; numbered last only to preserve the existing RQ1--RQ7 cross-references.)
\end{enumerate}

\subsection{RQ1: Does the pre-2022 anchor license a bounded, non-uniform false-positive guarantee at all, as opposed to a flat one or none?}

\textbf{Setup.} On the primary single-source pilot corpus (arXiv abstracts, $\approx$3,000 documents per year-bucket, thresholds frozen on the 2021 calibration split per Assumption A4), the pooled false-positive rate across all pre-2022 buckets is computed per detector family, together with its 95\% Clopper--Pearson interval, and the per-bucket upper confidence bound on false-positive rate is computed for every year.

\textbf{Expected outcome.} Both hold for at least 2 of the 3 detector families: (i) the pooled anchor false-positive rate is nonzero --- its 95\% Clopper--Pearson lower bound exceeds 0.25\%; and (ii) the bound is not uniform in time --- the maximum per-year upper confidence bound exceeds twice the nominal calibrated rate.

\textbf{Falsification condition.} Refuted if either conjunct fails for a majority of detector families: if the pooled lower bound is indistinguishable from zero, there is no anchor-erosion phenomenon to attribute at all; if the per-year upper confidence bound never exceeds twice the nominal rate, the anchor's guarantee is already uniform across the pre-2022 span and every downstream disentanglement question (RQ2--RQ7) becomes moot rather than merely smaller in magnitude.

\textbf{Hypothesis (plain language).} The pre-2022 anchor's implied error guarantee is real (not simply zero noise) and it is not flat across the years the anchor spans --- establishing that there is a genuine bound-shaped phenomenon for the rest of this study to attribute, propagate, and envelope, rather than presupposing one.

\wj{Standpoint (pre-registered H1): we expect BOTH conjuncts to hold for $\geq$2/3 families --- pooled 95\% Clopper--Pearson lower bound on FPR $>0.25\%$ (the anchor's guarantee is not zero-error) AND $\max_y$ upper bound $>2\nu=2\%$ (not uniform in time). If either fails for a majority of families, the erosion premise collapses and the paper pivots to the pre-registered robustness-certification framing (global fallback): the pre-2022 anchor is certified flat within the tested envelope --- still a publishable, decision-relevant finding.}

\textbf{Metrics.} Pooled false-positive rate and its Clopper--Pearson interval per detector family; per-bucket upper confidence bound on false-positive rate; the validity envelope $E_k(\epsilon,\delta)$ definition (Problem Definition) evaluated at the pre-registered tolerance.

\textbf{Result.} \placeholder{fill in after experiments are run}

\subsection{RQ2: Does the false-positive-rate drift survive restriction to a leakage-immune, pre-LLM sub-range of years?}

 \textbf{Setup.} Using the primary single-source pilot corpus (arXiv abstracts, ~3,000 documents per year-bucket, with detector thresholds frozen on a 2021 calibration split disjoint from all evaluation years), the drift regression is refit restricted to years at or before 2018 only, for each of three instance-level detector families.

\textbf{Expected outcome.} The slope remains significant and positive --- the same sign as the full-range slope --- for at least 2 of the 3 detector families when restricted to the leakage-immune sub-range.

\textbf{Falsification condition.} Refuted if the slope becomes non-significant or flips sign for a majority of detector families once restricted to years at or before 2018 --- that would mean the full-range slope was being driven by years adjacent to the 2022 cutoff, consistent with undetected leakage rather than a leakage-immune style effect.

\textbf{Hypothesis (plain language).} Whatever drift is observed over the full pre-2022 range is not simply an artifact of years sitting close to the cutoff, where leakage risk is highest.

\wj{Standpoint (pre-registered, part of H2): we expect the positive slope to survive on the $\leq$2018 sub-range (and the $\leq$2016 sensitivity cut to agree in sign) for $\geq$2/3 families --- i.e., the drift is NOT an artifact of cutoff-adjacent years where leakage risk concentrates. If it fails, the full-range slope is attributed to possible leakage and the stylistic-drift claim (and only that claim) is withdrawn under the pre-committed downgrade wording.}

\textbf{Metrics.} Per-bucket false-positive rate and confidence interval restricted to the sub-range; regression slope and confidence interval on the restricted range.

\textbf{Result.} \placeholder{fill in after experiments are run}

\subsection{RQ3: Does synthetic stylistic perturbation alone reproduce the direction of the observed drift?}

 \textbf{Setup.} A single fixed year's human sample is perturbed at five graded intensities of lexical and syntactic change (synonym substitution rate, sentence-length jitter, punctuation-convention shifts), and each detector's flag rate is measured at each intensity level, using the same fixed sample across all intensities (a paired design).

\textbf{Expected outcome.} Flag rate increases monotonically with perturbation intensity for at least 2 of the 3 detector families, replicating the sign of the slope observed on real year-to-year data, with an absolute flag-rate increase of at least 1 percentage point between the lowest and highest perturbation intensity (a pre-registered effect-size floor --- significance without magnitude does not pass).

\textbf{Falsification condition.} Refuted if flag rate is flat or non-monotonic across perturbation intensity for a majority of detector families, or if the absolute increase between the lowest and highest intensity falls short of 1 percentage point even where a monotone trend is statistically detectable --- either way, detectors are insensitive (or only negligibly sensitive) to the kind of stylistic change the drift hypothesis requires, undermining the claim that real year-over-year drift reflects a style mechanism at all.

\textbf{Hypothesis (plain language).} Stylistic change of the kind that plausibly occurs over a decade is, by itself, sufficient to move these detectors' flag rate --- establishing that a style mechanism is at least available, independent of whether it is what actually happened historically.

\wj{Standpoint (pre-registered H2b, with the 1-pp effect floor): we expect a monotone flag-rate rise across the five perturbation intensities with $\geq$1~pp absolute increase lowest-to-highest for $\geq$2/3 families --- establishing the style mechanism is AVAILABLE at a practically relevant magnitude, independent of what happened historically. Statistically detectable but $<$1~pp does not pass: significance without magnitude is pre-registered as a fail.}

\textbf{Metrics.} Flag rate at each perturbation intensity; trend statistic and significance on flag rate versus intensity.

\textbf{Result.} \placeholder{fill in after experiments are run}

\subsection{RQ4: How much residual impurity does the certified-human anchor actually carry?}

 \textbf{Setup.} A human-audited spot-check of $n=2{,}000$ documents, stratified proportionally across year buckets and sources, is reviewed for misdating, later heavy revision, translation, or ghostwriting --- the sample size powered against the envelope's 1\% impurity tolerance (at a plausible true rate of 0.3\%, the one-sided exact-binomial 95\% upper bound is $\approx$0.7\% < 1\%, so the audit can demonstrate compliance rather than merely fail to reject). Per Assumption A1, the audit runs as a validated instrument, not a pass-through metadata read: the rubric is piloted on a ~100-document slice, two independent annotators label the same sample, inter-rater agreement of $\kappa \geq 0.7$ is required before the impurity estimate is trusted (if $\kappa < 0.7$, the rubric is refined and re-validated before any impurity number is reported), and disagreements are resolved by a third annotator's majority vote.

\textbf{Expected outcome.} The estimated impurity rate stays below 1\%, with its upper confidence bound smaller than the magnitude of any detected drift effect from RQ2/RQ3.

\textbf{Falsification condition.} Refuted if the estimated impurity rate's upper confidence bound exceeds 1\% or is comparable to or larger than the detected drift magnitude --- in that case, any observed false-positive-rate slope cannot be safely attributed to style drift rather than to genuine mislabeled contamination, and every downstream claim in this study would need to be reported net of that uncertainty rather than as a clean drift estimate.

\textbf{Hypothesis (plain language).} The anchor's own authorship certification, while not perfect, is clean enough that its residual uncertainty does not swallow whatever drift effect the study is trying to measure.

\wj{Standpoint (pre-registered assumption gate, not a drift hypothesis): we expect estimated impurity $<1\%$ with its one-sided 95\% upper bound below the detected drift magnitude ($n=2{,}000$ is powered so a plausible 0.3\% true rate demonstrates compliance, upper bound $\approx$0.7\%). If the bound exceeds 1\%, every downstream bound is reported net of impurity and the drift claim must additionally survive the $\leq$2018 arm --- the audit can weaken conclusions but is designed so it can never silently strengthen them.}

\textbf{Metrics.} Annotator-flagged impurity rate and its one-sided 95\% upper confidence bound; inter-annotator agreement (Cohen's $\kappa$, acceptance gate $\kappa \geq 0.7$); the widening the bound contributes to each year-bucket's FPR confidence interval before the envelope is drawn (the audit's only channel into downstream results).

\textbf{Result.} \placeholder{fill in after experiments are run}

\subsection{RQ5: Does the false-positive-rate drift survive composition-matched resampling?}

 \textbf{Setup.} Each year bucket in the primary single-source corpus is resampled to equalize topic-cluster and document-length distribution against a reference bucket, at matched per-bucket token budgets, and the drift regression is refit on this composition-matched resample.

\textbf{Expected outcome.} The slope survives --- same sign, comparable magnitude, still significant at the pre-registered threshold --- in the composition-matched arm.

\textbf{Falsification condition.} Refuted if the slope disappears or reverses sign once composition is matched --- that would mean the originally observed trend was a genre or topic-composition artifact of corpus assembly, not a temporal style effect, and the stylistic-drift claim specifically would need to be withdrawn even though the anchor's overall invalidity could still stand under a different mechanism.

\textbf{Hypothesis (plain language).} The observed relationship between year and false-positive rate is a property of time itself, not a byproduct of which topics or genres happen to dominate the corpus's older versus newer years.

\wj{Standpoint (pre-registered, part of H2): we expect the slope to survive --- same sign, comparable magnitude, $p<0.0167$ --- after equalizing topic, length, genre/register, and source across buckets. If composition matching kills the slope, the trend was a corpus-assembly artifact; the anchor is still unsafe as conventionally used, but the mechanism claim changes under pre-committed wording (composition, not style).}

\textbf{Metrics.} Matched-arm slope and confidence interval; a per-bucket composition-similarity diagnostic relative to the reference bucket.

\textbf{Result.} \placeholder{fill in after experiments are run}

\subsection{RQ6: Does anchor age bias the corpus-level prevalence estimate a downstream user actually consumes?}

 \textbf{Setup.} A target corpus is held fixed --- both a synthetic mixture of known AI-fraction built from anchor human text plus an AI-side reference at a known mixing ratio, and a real post-2022 corpus of unknown ground truth --- while the human reference bucket fed into the mixture-based prevalence estimator is swapped between the most recent tested bucket and the most temporally distant tested bucket.

\textbf{Expected outcome.} The estimated AI-fraction shifts upward by at least 2 percentage points when the aged reference is substituted, holding the target fixed, significant under a paired permutation test.

\textbf{Falsification condition.} Refuted if the estimated AI-fraction is invariant, within its confidence interval, to which reference bucket is used --- meaning the mixture estimator is robust to anchor age, and any per-document drift detected in RQ2/RQ3/RQ5/RQ8 does not propagate to the corpus-level number end users actually consume. \textbf{Positive branch (pre-registered, co-headline):} if RQ2/RQ5/RQ8 establish drift but this question resolves negative, the finding is reported in its positive form --- the estimator is self-robust to anchor age: per-document false-positive rates drift measurably with temporal distance, yet the corpus-level prevalence estimate is invariant to reference age within 2 pp, so the pre-2022 convention is safe at the corpus level even though its documents drift. \textbf{Inconclusive band:} if the permutation confidence interval on the shift straddles the 2 pp floor with half-width above 1 pp, the propagation question is reported INCONCLUSIVE (widen the target corpus or add reference buckets), never collapsed into either branch.

\textbf{Hypothesis (plain language).} Anchor age does not only move an intermediate detector statistic --- it moves the actual prevalence number a downstream user would report and act on.

\wj{Standpoint (pre-registered H3, CO-PRIMARY claim): we expect the aged reference to shift the estimated AI-fraction upward by $\geq$2~pp (paired permutation $p<0.05$). BOTH branches are pre-registered as publishable: propagation $\geq$2~pp indicts the convention where users feel it; invariance $<$2~pp is reported as the positive co-headline --- the estimator is self-robust to anchor age even though documents drift. Inconclusive band: permutation CI straddling the 2-pp floor with half-width $>$1~pp reports INCONCLUSIVE, never either branch.}

\textbf{Metrics.} The shift in estimated AI-fraction between reference-bucket conditions; permutation-test significance; on the known-composition target, the absolute bias against the true mixing ratio.

\textbf{Result.} \placeholder{fill in after experiments are run}

\subsection{RQ7: Does a confidence-bound validity envelope draw a meaningfully different boundary than a naive significance-threshold envelope?}

 \textbf{Setup.} On the same fitted drift curve, the validity envelope is constructed two ways --- a naive construction that draws the boundary at the first year the raw slope's p-value crosses significance, and a confidence-bound construction that draws the boundary at the largest contiguous range whose upper confidence bound on false-positive rate stays within a pre-declared tolerance of the nominal calibrated rate --- and the two boundaries are compared directly.

\textbf{Expected outcome.} The confidence-bound envelope is strictly narrower --- draws the safe boundary at a more recent year --- than the naive envelope, for at least one tested source or genre.

\textbf{Falsification condition.} Refuted if the two constructions draw an identical boundary across all tested sources and genres --- that would mean the confidence-bound construction does no additional work relative to the simpler significance-threshold rule, and the more conservative bookkeeping would be unnecessary complexity rather than a load-bearing contribution.

\textbf{Hypothesis (plain language).} A boundary drawn only from statistical significance is not the same thing as a boundary drawn from genuine certified safety, and the difference between the two is large enough to matter in practice.

\wj{Standpoint (pre-registered, envelope-construction check): we expect the confidence-bound envelope to draw a strictly different (more recent, more conservative) safe boundary than the naive significance-threshold rule for at least one source/genre --- demonstrating the construction is load-bearing where old buckets thin out. If the two boundaries coincide everywhere, we report the simpler rule as sufficient and the construction as redundant --- an honest negative on our own machinery.}

\textbf{Metrics.} The boundary year under each construction, per source/genre; the gap, in years, between the two boundaries.

\textbf{Result.} \placeholder{fill in after experiments are run}

\subsection{RQ8: Does the false-positive-rate drift survive a provenance-stratified, pipeline-constant arm?}

 \textbf{Setup.} Within at least one corpus, the digitization pipeline is held constant across all years --- born-digital text only, single source, single processing pipeline (e.g., arXiv abstracts retrieved from the arXiv's own metadata API, or a single venue's own review archive) --- and the drift regression is refit on this arm, at matched per-bucket n and token budgets. Attribution rank: alongside RQ2 (leakage-immune) and RQ5 (composition-matched) --- the third one-alternative-immune control, targeting the processing-history rival that neither of the others covers.

\textbf{Expected outcome.} The slope survives --- same sign, significant at the pre-registered threshold --- in the pipeline-constant arm for at least 2 of the 3 detector families.

\textbf{Falsification condition.} Refuted if the slope vanishes or reverses sign in the pipeline-constant arm for a majority of detector families --- the drift is then re-attributed to processing artifacts, with the pre-committed wording: the pre-2022 anchor's apparent temporal drift is an artifact of heterogeneous digitization pipelines, not of writing-style change; the anchor is unsafe as conventionally assembled from mixed-provenance pools but safe when restricted to pipeline-constant sources, making the validity envelope a provenance condition rather than a time condition. (This branch re-scopes Claim 3 and drops Claim 2 only, exactly as the composition-matched downgrade does.)

\textbf{Hypothesis (plain language).} The observed drift reflects how people wrote, not how their text was later scanned, encoded, or normalized.

\wj{Standpoint (pre-registered, part of the attribution family): we expect the slope to survive in the pipeline-constant (born-digital, single-source, single-pipeline) arm for $\geq$2/3 families --- the drift reflects how people wrote, not how their text was later scanned or normalized. If it vanishes, the pre-committed re-attribution applies: the anchor is unsafe as conventionally assembled from mixed-provenance pools but safe when pipeline-constant, making the envelope a provenance condition rather than a time condition --- still decision-relevant.}

\textbf{Metrics.} Pipeline-constant-arm slope and confidence interval per detector family; a pipeline-homogeneity diagnostic (encoding, markup, and OCR-artifact incidence per bucket) confirming the arm is actually constant.

\textbf{Result.} \placeholder{fill in after experiments are run}


\subsection{Further Evaluation}
With more time, a natural extension is broader language and genre coverage beyond the sources tested here, since the validity envelope this design produces is explicitly source- and genre-scoped rather than a universal claim. A second natural extension, deliberately out of scope for this paper, is proposing a drift-robust prevalence estimator that corrects for the measured bias rather than only reporting it --- that is a repair, not a measurement, and belongs to a different contribution than the one made here.

\section{Threats to Validity}\label{sec:threats}

\textbf{Construct validity.} The ``certified-human'' label is a publication-year timestamp, a proxy for authorship rather than a verified property: ghostwriting, later heavy revision, translation, or misdating can all survive the metadata filter. We do not assume this away --- a frozen human audit ($n=2{,}000$ documents stratified across year-buckets and sources, two independent annotators, $\kappa \geq 0.7$ acceptance gate, third-annotator adjudication) bounds the residual impurity, and that bound widens every per-bucket false-positive confidence interval \emph{before} the envelope is drawn, so label uncertainty can render a bucket inconclusive but can never manufacture a spurious ``safe.''

\textbf{Internal validity.} Calendar year is entangled with three rivals to stylistic drift. Corpus composition (topic, length, genre/register, source) is addressed by the composition-matched resampling arm, which equalizes all four named covariates across buckets. Leakage of AI text into late pre-2022 buckets is addressed by the leakage-immune sub-range: the boundary is set at 2018 because GPT-2's February 2019 release marks the first widely available fluent open-domain generator, and a $\leq$2016 sensitivity cut must agree in sign. The processing-history confound --- OCR error profiles, encoding artifacts, and normalization differences that track the year of \emph{digitization} rather than authorship --- is addressed by the provenance-stratified pipeline-constant arm (born-digital, single source, single pipeline); if the slope vanishes there, the drift is re-attributed to processing artifacts under pre-committed wording rather than silently retained. The instrument itself is guarded in both directions: a year-shuffle placebo must report a null slope (the pipeline cannot invent a trend) and a zero-distance sanity cell must report near-nominal false-positive rates (it does not mask one).

\textbf{External validity.} The validity envelope is explicitly scoped to the tested sources, genres, and languages --- English academic abstracts, peer reviews, news, books, and forum text --- and is reported as a source- and genre-conditional bound, never a universal claim. Detector families released after this study, and non-English corpora, are outside the envelope by construction; broadening that coverage is the natural replication path.

\textbf{Conclusion validity.} Old year-buckets are thin, and an under-powered bucket must not be confused with a safe one. A dual trigger operationalizes this: a bucket enters the envelope only if $n_y \geq n_{\min} = 1{,}000$ and its realized Clopper--Pearson half-width is $\leq 1$~pp; otherwise it reports inconclusive. All detector-level tests are Bonferroni-corrected across the three pre-registered families, effect-size floors accompany every significance threshold (significance without magnitude does not pass), and every hypothesis carries a pre-committed fallback wording, so no branch of the analysis requires post-hoc reframing.

\section{Related Work}\label{sec:related}

The closest prior work to this paper is the corpus-level mixture-estimation framework that popularized the pre-2022 anchor \cite{MonitorAIScale}, and it is worth stating the delta plainly before surveying the wider landscape: that framework fits a single mixing coefficient against a fixed human reference distribution and treats it as unproblematic ground truth, but it never tests whether the reference's temporal composition confounds the estimated AI fraction --- which is precisely the instrument this paper audits directly, rather than merely reusing. The absence claim underlying this positioning is auditable, not asserted: a systematic search over 11,683 OpenReview records was run from five angles --- (a) AI/machine-generated-text detection and detector false-positive behavior; (b) corpus-level prevalence/proportion estimation of AI content; (c) temporal distribution shift, concept drift, and temporal-holdout evaluation; (d) benchmark/data-contamination detection and its assumption fragility; (e) pre-training-data/membership-inference detectors and calibration --- and returned no work that measures a pre-2022 certified-human anchor's temporal validity.

\textbf{Corpus-level AI-content prevalence estimation.} The mixture-estimation framework \cite{MonitorAIScale} reframes AI-text detection from instance classification to a population-level measurement problem, fitting a single mixing coefficient between a human-written reference distribution and an AI-generated reference distribution over a stable vocabulary, demonstrated on AI-conference peer reviews. This establishes the corpus-level framing this paper's Results depend on, but treats the human reference corpus as a fixed, unproblematic input rather than an object whose own temporal validity might need auditing --- exactly the gap this paper's reference-age propagation test is built to close. Its own Limitations section explicitly names ``the temporal distribution shift in token frequencies'' as an acknowledged but unmeasured source of error, which this paper turns into the central measured quantity rather than a disclaimed afterthought.

\textbf{Temporal concept drift and longitudinal benchmark validation.} A longitudinal malware-detection benchmark spanning more than a decade \cite{LAMDA} holds a consistent feature space fixed across years specifically to measure model degradation over time cleanly, motivated by evidence that naive temporal evaluation inflates performance estimates. This work supplies the year-stratified, temporally-structured validation template this paper adapts, but it targets classification-accuracy decay in a fixed feature space for a different domain (malware), not a calibrated false-positive-rate bound for a text negative-anchor --- its protocol is borrowed, not its findings. A temporal-cutoff contamination diagnostic in financial forecasting \cite{DetectingLookaheadBiasIn} is structurally analogous to this paper's own year-shuffle placebo and pre-LLM sub-range disentanglement --- an interaction test that must vanish in a post-cutoff placebo sample --- confirming the same design pattern generalizes across domains even though its target task is unrelated to corpus-level AI-content detection.

\textbf{Contamination- and membership-detector assumption fragility.} Work stress-testing benchmark-contamination detectors \cite{BenchContamFragility} shows that the assumption underlying most such detectors --- that seen and unseen samples remain statistically distinguishable --- breaks down under staged post-training, across detector families spanning reference-based, reference-free, perturbation-based, and output-distribution approaches. A reference-free pre-training-data detector \cite{MinKpp} shows a closely related failure mode: membership heuristics degrade toward random as the training and non-training distributions being compared converge. A collection-level contamination detector \cite{ContamICL} further motivates measuring at the corpus level specifically because sample-level signals are confounded by distribution shift. A calibration layer for metric-based detectors \cite{MarkovCalib} shows that detector reliability is brittle under paraphrasing and distribution shift at the mechanism level. All four diagnose that detector reliability collapses as compared distributions converge or shift, but none turns that diagnosis into an explicit false-positive-rate bound as a function of temporal distance for a certified-human anchor --- the temporal gap this paper measures is an unaudited special case of exactly the phenomenon this literature describes only in the abstract.

\textbf{On the envelope construction itself.} We concede that the confidence-bound envelope uses no new statistical machinery --- the per-year upper bounds are standard Clopper--Pearson intervals, and the change-point fit is ordinary segmented regression. The contribution is not the estimator but the *estimand* it is applied to: the anchor-conditional false-positive rate as a function of temporal distance from the certification cutoff, a quantity no prior work defines or reports. Using textbook intervals to bound a previously-unmeasured quantity is the intended shape of an analysis paper, and RQ7 tests directly whether the confidence-bound envelope draws a materially different safe boundary than the naive significance-threshold rule it replaces, so the construction earns its place empirically rather than by assumption.

\textbf{Per-document detection instruments.} A zero-shot detector \cite{DetectGPT} thresholds the local probability-curvature of a passage under a source model, needing no trained classifier or labeled data; a rewrite-based detector \cite{LearnToDistance} instead learns the distance between a passage and its LLM rewrite via normalized log-likelihood rather than a fixed edit or semantic distance; a single-call black-box detector \cite{DnR} shuffles tokens within chunks and uses recovery similarity as a detection proxy, suited to closed-model, low-query deployment. Each of these is validated against contemporaneous human and AI samples, and each is used in this paper's own estimator battery as an instrument --- not a baseline to be beaten, since this is a measurement study, but a component whose false-positive rate on temporally distant human text is exactly what remains uncharacterized in its own original evaluation. A held-out benchmark spanning many models, domains, and adversarial attacks \cite{RAID} is used only to validate that each battery member reproduces its expected behavior before any threshold is frozen, per Assumption A4.

\textbf{Structural exemplars for assumption-auditing measurement studies.} Two recent empirical-study exemplars shape this paper's evidence structure rather than its content: a controlled study interrogating whether larger pretrained language models are actually better calibrated \cite{ACloseLookInto}, and a study measuring whether model-compression techniques degrade fairness across multiple methods and metrics \cite{AComparativeStudyOn}. Both fall on the correct side of the archetype for this paper --- a field-wide operating assumption is turned into a falsifiable, multi-condition measurement rather than argued rhetorically --- and both are imitated here for evidence shape rather than for content.

\textbf{An alternative, fix-oriented response to temporal drift.} A contemporaneous line of work \cite{MovingTargetTTA} responds to the same underlying observation --- that detector performance degrades under continual distribution shift, including drift in human writing style --- not by measuring and bounding the drift, but by proposing test-time adaptation: an online, semi-supervised procedure that re-calibrates a detector against unlabeled data as it arrives. This is a genuinely different contribution class from the one made here: a repair method with a method-superiority evidence standard, rather than a measurement study with a validity-envelope evidence standard. The two are complementary rather than competing --- a validated envelope tells a practitioner when reliance on a frozen anchor becomes unsafe; test-time adaptation is one candidate response once it does.

**

\section{Conclusion}\label{sec:conclusion}

This work reframes a widely-used calibration convention --- that pre-2022 human text is a self-certifying, timeless negative anchor for corpus-level AI-content estimators --- as an empirical question with a precise attribution structure, rather than an assumption to take on faith. The central methodological wager is that a false-positive-rate drift is only informative once it survives controls that each rule out a specific rival explanation (composition confound, content leakage, estimator staleness) and its practical stakes are only established once it is shown, not assumed, to move the corpus-level number a downstream estimator actually reports. Whichever outcome the measurement yields, the finding is decision-relevant: a genuine, propagating drift would mean every corpus-level AI-prevalence estimate that reuses this anchor without stating a validity envelope carries an uncharacterized, directional bias that grows as the anchor ages further past 2022; a flat, non-propagating result would instead hand the field its first quantified license for a convention it has so far only assumed. Two limitations bound this claim from the outset and are natural next steps rather than afterthoughts: no scalable oracle can fully certify that a nominally pre-2022 document involved no later revision or mislabeling, so the reported bounds are conditional on the purity spot-check's own coverage, not an absolute guarantee; and the validity envelope this design produces is scoped to the tested sources, genres, and languages, leaving broader-coverage replication --- and any eventual drift-robust estimator repair --- for follow-up work.

---

\section*{Data Availability}
The year-stratified certified-human anchor corpus (bucket manifests with hard archival timestamps and provenance metadata), the detector-battery and drift-curve estimation code, the composition-matching and pipeline-constant resampling scripts, the impurity-audit rubric with anonymized annotator labels, and the full pre-registration document (hypotheses, thresholds, effect-size floors, fallback wordings, and analysis gates, all frozen before data collection) will be released as a replication package upon publication. \placeholder{Artifact DOI/URL to be inserted at camera-ready.}

\bibliographystyle{IEEEtranN}
\bibliography{refs}

\end{document}
